%%%%%%%%%%%%%%%%%%%%%%%%%%%%%%%%%%%%%%%%%%%%%%%%%%%%%%%%%%%%%%%%%%%%%%%%%%%%%%%%
% Obxectivo: Lista de termos empregados no documento,                          %
%            xunto cos seus respectivos significados.                          %
%%%%%%%%%%%%%%%%%%%%%%%%%%%%%%%%%%%%%%%%%%%%%%%%%%%%%%%%%%%%%%%%%%%%%%%%%%%%%%%%

\newglossaryentry{bytecode}{
  name=bytecode,
  description={Código independente da máquina que xeran compiladores de determinadas linguaxes (Java, Erlang,\dots) e que é executado polo correspondente intérprete.}
}

\newglossaryentry{Escena}{
  name=Escena,
  description={Las Escenas son el tipo de archivo básico con el que se trabaja en Godot. Estas son una estructura jerárquica de \gls{Nodo}s, que parten de un único Nodo raiz. No confundir con \gls{escena} }
}

\newglossaryentry{Nodo}{
  name=Nodo,
  description={Los Nodos son la estructura lógica mínima con la que se puede trabajar en Godot. Todo elemento que vaya a estar presente en tiempo de ejecución en el programa es un Nodo. Funcionan como las clases en la programación orientada a objetos, puesto que cada tipo de Nodo cuenta con atrubutos, funciones y señales propias, a parte de heredar las su clase padre. No confundir con \gls{nodo}}
}

\newglossaryentry{escena}{
  name=escena,
  description={Espacio visual en la pantalla. Sinónimo de ``vista''. No confundir con \gls{Escena}}
}

\newglossaryentry{nodo}{
  name=nodo,
  description={Punto de origen de una ramificación. En el contexto de árboles de decisión: representación de una bifurcación producto de una decisión del árbol basada en un atributo de los datos de entrenamiento. En el contexto de redes neuronales: representación del producto escalar entre el vector de salida de los nodos previos y el vector de pesos propio más un valor de sesgo; sinónimo de ``neurona''. No confundir con \gls{Nodo}}
}
